% ============================================================================
\chapter{IRQ、PIC 与键盘驱动}
\label{ch:irq-keyboard}
% Stage 3 — 对应 2026-03-10 changelog（IRQ + Keyboard）
% ============================================================================

\section{本章目标}

\begin{itemize}
  \item 理解 8259 PIC（可编程中断控制器）的工作原理
  \item Remap PIC 将 IRQ0--15 映射到 IDT 向量 0x20--0x2F
  \item 实现 IRQ 分发框架（注册 handler + EOI）
  \item 实现 PS/2 键盘驱动：scancode → ASCII + 环形缓冲区
\end{itemize}

\section{背景知识：8259 PIC}

% TODO:
% - 主 PIC (IRQ0-7, port 0x20/0x21) + 从 PIC (IRQ8-15, port 0xA0/0xA1)
% - ICW1-ICW4 初始化序列
% - Remap: 为什么要从默认的 0x08 改到 0x20（避开 CPU 异常）
% - IMR (Interrupt Mask Register): 屏蔽/使能 IRQ
% - EOI (End of Interrupt): 告诉 PIC 中断处理完毕

\subsection{IRQ 向量映射}

% TODO: 表格 IRQ0=0x20(Timer), IRQ1=0x21(Keyboard), ...

\section{PIC 初始化：\texttt{pic.c}}

% TODO: pic_init() 的 ICW1-4 序列逐步解释

\section{IRQ 分发框架：\texttt{irq.c}}

% TODO: g_irq_handlers[16], irq_install_handler, irq_handler_c

\section{IRQ 汇编 Stub：\texttt{irq\_stubs.asm}}

% TODO: 和 ISR stub 类似，但 push 的是 IRQ 号 + 调用 irq_handler_c

\section{PS/2 键盘驱动：\texttt{keyboard.c}}

% TODO:
% - 端口 0x60 读 scancode
% - Scancode Set 1 → ASCII 转换表（含 Shift 状态）
% - 环形缓冲区（ring buffer）存储按键
% - keyboard_getc() / keyboard_try_getc()

\begin{designbox}
为什么用环形缓冲区？中断处理程序必须尽快返回（否则会阻塞其他中断）。
把字符存入缓冲区、由 shell 主循环消费，是经典的"生产者-消费者"模式。
\end{designbox}

\section{验证}

\begin{lstlisting}[style=shellstyle]
make iso && make run
# 在 QEMU 窗口按键，串口应看到按键回显
\end{lstlisting}

\section{本章小结}

硬件中断链路打通：PIC → IRQ stub → C handler → 键盘驱动 → 缓冲区。
下一章我们用这个输入做交互式 Shell。
